\section{Related Work}
\label{sec:related}

\subsection{Spine MRI Grading}
% Windsor SpineNetV2, RSNA top-10 + published RSNA 2024 method papers,
% SpineNet external validations, SPIDER benchmark. Position our
% contribution: 3D volumetric grading with cross-dataset zero-shot
% label extension and the first BiomedCLIP-based RSNA 2024 pipeline.
SpineNetV2~\cite{windsor2024spinenet} provides a 3D ResNet-34 backbone
pretrained on lumbar IVD volumes, which we adopt as our base. Several
external-validation studies place SpineNet's per-pathology agreement at
$\kappa$~0.46--0.74 across Pfirrmann, central canal stenosis,
spondylolisthesis, foraminal stenosis, and disc herniation
\cite{nigru2024spinenetv2,mcsweeney2023spinenet,liu2025spinenetv2ext},
with foraminal grading consistently the weakest sub-task ($\sim70\,\%$
balanced accuracy ceiling). The RSNA~2024 challenge produced multiple
multi-stage 2D pipelines~\cite{nshen7rsna2024,kaggle2nd2024,kaggle4th2024};
these treat each IVD as a 2D crop, fuse sagittal and axial planes via
ensembling, and report only the competition's sample-weighted log-loss
metric. The single peer-reviewed RSNA~2024 method paper to date,
Aktan et al.~\cite{liu2025rsna}, evaluates a MobileNetV2-UNet
joint segmentation+classification cascade on a balanced internal split.
Other recent multi-task or per-condition baselines include
Hallinan et al.~\cite{hallinan2021stenosis} (seminal Mask R-CNN + 3D CNN
stenosis classifier with dichotomous $\kappa$~0.89--0.96),
Lin et al.~\cite{lin2024multiattention} (the closest CBAM-based
architectural twin, evaluated on a binarised single-level L4--L5 subset),
M-SCAN~\cite{hong2025mscan} (multi-view cross-attention on RSNA 2024,
AUROC~0.971 on canal stenosis), and the
Phaphuangwittayakul et al.~\cite{phaphuangwittayakul2026multitask}
multi-task model with multi-center external validation. Our approach
differs in three ways: (i) we use volumetric input throughout (no 2D
slice ensembling); (ii) we are, to our knowledge, the first to integrate
a frozen vision-language foundation model on RSNA~2024 lumbar grading;
and (iii) we add a label-space-extensible head allowing inference on
unseen disease classes from the SPIDER dataset~\cite{vandergraaf2024spider,spider}
without any SPIDER training.

\subsection{Attention for Imbalanced Medical Imaging}
% CBAM (Woo 2018), SE-Net (Hu 2018), small-lesion attention precedents.
CBAM~\cite{woo2018cbam} combines channel and spatial attention at near-zero
parameter cost. We choose CBAM over SE-Net~\cite{hu2018senet} because the
spatial gate is the critical component for localising small Severe
lesions; non-local~\cite{wang2018nonlocal} and transformer self-attention
scale poorly under our small-data regime.

\subsection{Medical Vision-Language Foundation Models}
% BiomedCLIP, MedGemma, RadFM, RAD-DINO, RadCLIP, Decipher-MR, Merlin.
% Justify why a frozen contrastive encoder is the right choice.
Recent medical foundation models include BiomedCLIP~\cite{zhang2023biomedclip},
RadCLIP~\cite{lu2024radclip}, RadFM~\cite{wu2025radfm},
RAD-DINO~\cite{perezgarcia2024raddino}, MedGemma~\cite{sebro2025medgemma},
the 3D-MRI-native Decipher-MR~\cite{yang2026deciphermr}, and the 3D-CT
Merlin~\cite{blankemeier2026merlin}. Of these, BiomedCLIP is the only
contrastive image--text encoder small enough (${\sim}86$M image params)
to integrate as a frozen feature extractor on a consumer GPU while
exposing a CLIP-style~\cite{radford2021clip} embedding space suitable
for cosine-similarity zero-shot classification. Recent analyses
\cite{biomedclipImbalance2025} document BiomedCLIP's brittleness under
high class imbalance, which motivates our hybrid design that pairs the
frozen 2D encoder with a learnable 3D-CBAM-ResNet specialist.
MedGemma and RadFM are generative models whose multimodal arms are
2D-only and whose fine-tuning requirements exceed our resource regime;
RadCLIP introduces a slice-pooling adapter that would natively handle
volumetric inputs but appeared after our experiments were finalised;
Decipher-MR and Merlin are very recent native-3D foundation models
which we discuss as the strongest future-work alternatives in
Sec.~\ref{sec:discussion}. To our knowledge, no published study has
applied BiomedCLIP (or any of these alternatives) to RSNA~2024 lumbar
spine MRI grading prior to ours.
